\chapter*{Conclusion}
\addcontentsline{toc}{chapter}{Conclusion}
\label{cap:conclusione}
In this thesis, we presented an algebraic reconstruction of classical and quantum mechanics along with their subsequent geometric unification. In the first chapter, we rebuilt classical phase space from a commutative $C^*$-algebra via Gelfand-Naimark, supplemented by a Poisson bracket to endow it with a symplectic structure. In the second chapter, we reconstructed quantum mechanics from a non-commutative algebra using the GNS construction to obtain Hilbert space representations. Hence, the Poisson bracket was replaced by a Lie bracket, establishing a formal correspondence with the classical framework. In the third chapter, we joined the algebraic and geometric formulations by identifying the GNS-derived pure state space with projective Hilbert space. Its intrinsic Kähler structure generates Schrödinger dynamics from its symplectic form. Moreover, it encodes Born-rule probabilities through its metric, while expressing symmetries as conserved charges via a geometric moment map.

From this work, we can conclude the following key facts. In algebraic classical mechanics, once the topological state space is fixed, it is needed a choice of the smooth subalgebra, as in Definition \ref{def:smooth-subalgebra}. Instead, in the algebraic quantum picture, we observe that through GNS Theorem \ref{thm:gns} we recover the Dirac-Von Neumann formalism and a Lie algebra using only operational structures. At the end, once the theories are rewritten geometrically, we can conclude the quantum algebraic axiomatization using a classical geometrical analogy because we showed that same geometrical objects represents the same physical processes as in the classical case.

Limitations of this framework discussed in Section \ref{sec:concl-limits} and Remark \ref{rem:purification-stratified-spaces} can be overcome by extending the formalism to infinite dimensional Hilbert spaces as in \cite{bratteliOperatorAlgebrasQuantum} and mixed states in a future work. This thesis can also lead to tackling the choice problem of the smooth structure in Chapter \ref{cap:cm-c-star}. In fact, one can use the quantum Kählerian algebra in Definition \ref{def:kahlerian-function} on the projective space to identify a proper smooth Lie algebra, using it as a choice of a Poisson algebra of smooth functions. Another future work can focus on finding an operational interpretation and a subsequent ontology for this new algebraic-geometric framework. 

